\documentclass[11pt,a4paper]{article}
\usepackage{amsmath,amssymb,amsthm}
\usepackage{listings}
\usepackage{xcolor}
\usepackage{geometry}
\geometry{margin=1in}

\lstset{
  language=Lean,
  basicstyle=\ttfamily\small,
  keywordstyle=\color{blue},
  commentstyle=\color{green!50!black},
  stringstyle=\color{red},
  breaklines=true,
  frame=single,
  numbers=left
}

\title{\textbf{Hodge Conjecture as a dm³ System (Classical First)}\\
A Lean-First Language Specification \\
(Algebraic–Geometric Pillar of the Unified Framework)}
\author{Pablo Nogueira Grossi \\ G6 LLC / AXLE Project}
\date{v1.0 — April 2026}

\begin{document}

\maketitle

\begin{abstract}
Classical Hodge Conjecture in the dm³\_hodge framework.
Smooth projective varieties over $\mathbb{C}$, Hodge cycles vs algebraic cycles.
C = cohomological compression, K = curvature of the Hodge bundle,
F = folding via degenerations, U = unfolding to algebraic cycles.
This mirrors the other dm³ pillars and leaves explicit hooks for mixed and absolute Hodge.
\end{abstract}

\section{dm³ Grammar (Lean-first)}

\subsection{Syntax}
\begin{lstlisting}
inductive Dm3Op
| C | K | F | U
\end{lstlisting}

\subsection{Classical Interpretation}
C = cohomological compression (projection to Hodge pieces), \\
K = curvature of the Hodge bundle, \\
F = folding via degenerations / limiting mixed Hodge structures, \\
U = unfolding to algebraic cycles.

\subsection{Composition}
\[
G = U \circ F \circ K \circ C.
\]

\section{The dm³\_hodge Category (Classical)}

\subsection{Objects}
\begin{lstlisting}
structure Dm3HodgeObject :=
  (X           : HodgeVariety)
  (contactForm : HodgeVariety → HodgeVariety)
  (classical   : Prop := True)
\end{lstlisting}

\subsection{Morphisms}
\begin{lstlisting}
structure Dm3HodgeMorph (A B : Dm3HodgeObject) :=
  (map               : A.X → B.X)
  (preserves_contact :
    ∀ x, map (A.contactForm x) = B.contactForm (map x))
\end{lstlisting}

\section{Hodge as a dm³ Object (Classical)}

\subsection{State Space}
\begin{lstlisting}
def Hodge_state : Type := HodgeVariety
\end{lstlisting}

\subsection{Contact Form}
\begin{lstlisting}
def Hodge_contact : Hodge_state → Hodge_state := hodgeContact
\end{lstlisting}

\subsection{Object Definition}
\begin{lstlisting}
def Hodge_dm3 : Dm3HodgeObject :=
{ X           := Hodge_X,
  contactForm := Hodge_contact,
  classical   := True }
\end{lstlisting}

\section{The Hodge → dm³\_hodge Embedding}

\subsection{Operator Decomposition}
\begin{lstlisting}
theorem Hodge_operatorDecomposition :
  ∀ X : Hodge_state,
    Hodge_step X = (G C_Hodge K_Hodge F_Hodge U_Hodge) X := ...
\end{lstlisting}

\subsection{Contact Preservation}
\begin{lstlisting}
axiom Hodge_preserves_contact :
  ∀ X : Hodge_state,
    Hodge_step (Hodge_contact X) = Hodge_contact (Hodge_step X)
\end{lstlisting}

\section{Remaining Analytic Axioms (Classical Hodge)}

\subsection{meanContraction\_hodge}
\begin{lstlisting}
axiom meanContraction_hodge :
  ∀ X : Hodge_state,
    (∫ _ : Hodge_state,
        Real.log (hodgeCurvature (Hodge_step X) / hodgeCurvature X) ∂μ) < 0
\end{lstlisting}

\subsection{lyapunovDescent\_hodge}
\begin{lstlisting}
axiom lyapunovDescent_hodge :
  ∀ X : Hodge_state,
    hodgeHeight (Hodge_step X) < hodgeHeight X
\end{lstlisting}

\subsection{hasStructuredCycle\_hodge}
\begin{lstlisting}
axiom hasStructuredCycle_hodge :
  ∃ A : Set Hodge_state, is_dm3_hodge_cycle A
\end{lstlisting}

\section{Coherence Bridge Synchronization}

\begin{lstlisting}[language=YAML]
- id: hodge_conjecture
  claim_level: formally_stated
  lean4_sorry_label: "ADMIT-D_hodge, ADMIT-E_hodge, ADMIT-F_hodge"
  lean4_sorry_count: 3
  axle_status: "Dm3Hodge.lean v1.0 — Classical Hodge"
\end{lstlisting}

\section{Extensions: Mixed and Absolute Hodge}

Mixed and absolute Hodge structures can be added by replacing
\texttt{HodgeVariety} with appropriate categories and extending the
contact form to variations of mixed or absolute Hodge structures.
The dm³ operator grammar remains unchanged.

\end{document}
